% GENERATED FILE. Edit canonical lessons, then run npm run generate:books.
% canonical-source-hash: fnv1a64:0176cd52de4f04e7
% canonical-lessons: GU-C328-puja, GU-C328-arti, GU-C328-prasad, GU-C328-pragtavvu, GU-C328-dhup

\chapter{Evening Prayers at the Temple}
\label{ch:gu-evening-prayers-at-the-temple}
\hlchaptermodality{328}

\begin{chapteropening}
\textbf{By the end of this chapter:} \emph{I can say what happens at a temple in the evening.}

Complete the last lesson of chapter 328: I can say what happens at a temple in the evening.
\end{chapteropening}

\section[\texorpdfstring{\gu{પૂજા} (pūjā)}{પૂજા (pūjā)}]{\gu{પૂજા} (pūjā) — worship}

\label{lesson:GU-C328-puja}



\begin{quote}
Before the new one: say the Gujarati for a custom, then the Gujarati for a holy day.
\end{quote}

\subsection*{You'll want to know: \gu{પૂજા}}
\textbf{\gu{પૂજા}} — \emph{pūjā} — \textquotedblleft{}worship\textquotedblright{}.

\begin{quote}
\gu{સવારે} \gu{દાદી} \gu{મંદિરમાં} \gu{પૂજા} \gu{કરે} \gu{છે}.
\end{quote}

\begin{quote}
\emph{savāre dādī mandirmā̃ pūjā kare che.}
\end{quote}

\begin{quote}
In the mornings my grandmother worships at the temple.
\end{quote}

A feminine noun.

\begin{cousinweb}
\textbf{\gu{પૂજા}} (\emph{pūjā}) is a tatsama, Sanskrit \emph{pūjā}, \textquotedblleft{}honour, worship\textquotedblright{} (CDIAL 8317). It is done with \textbf{\gu{કરવું}}: \textbf{\gu{પૂજા} \gu{કરે} \gu{છે}}.
\end{cousinweb}

\begin{grammarlens}[title={What you've built}]
One more word for telling what happened.
\end{grammarlens}

\subsection*{Your turn}
\begin{itemize}
\raggedright
  \item \emph{Say it:} \emph{pūjā}
  \item \emph{Say it:} \emph{pūjā}, once more
  \item \emph{Say it:} \emph{savāre dādī mandirmā̃ pūjā kare che.} — In the mornings my grandmother worships at the temple.
  \item \emph{Recall:} say the Gujarati for a custom, then the Gujarati for a holy day, then say \emph{pūjā} again
\end{itemize}

\begin{tcolorbox}[breakable,colback=teal!4,colframe=teal!35!black,title={Before you move on}]
What does \gu{પૂજા} mean? (\textquotedblleft{}Worship\textquotedblright{}.) Say it once more.
\end{tcolorbox}
\section[\texorpdfstring{\gu{આરતી} (ārtī)}{આરતી (ārtī)}]{\gu{આરતી} (ārtī) — an aarti (worship with lamps)}

\label{lesson:GU-C328-arti}



\begin{quote}
Before the new one: say the Gujarati for a holy day, then the Gujarati for worship.
\end{quote}

\subsection*{You'll want to know: \gu{આરતી}}
\textbf{\gu{આરતી}} — \emph{ārtī} — \textquotedblleft{}an aarti (worship with lamps)\textquotedblright{}.

\begin{quote}
\gu{સાંજે} \gu{મંદિરમાં} \gu{આરતી} \gu{થઈ}.
\end{quote}

\begin{quote}
\emph{sā̃je mandirmā̃ ārtī thaī.}
\end{quote}

\begin{quote}
In the evening the aarti was held at the temple.
\end{quote}

A feminine noun.

\begin{cousinweb}
\textbf{\gu{આરતી}} (\emph{ārtī}) goes back to Sanskrit \emph{ārātrika}, the waving of lamps before an image, and CDIAL lists the Gujarati noun, \textquotedblleft{}the ceremony\textquotedblright{} (CDIAL 1315). It is feminine: \textbf{\gu{આરતી} \gu{થઈ}}, the aarti took place.
\end{cousinweb}

\begin{grammarlens}[title={What you've built}]
One more word for telling what happened.
\end{grammarlens}

\subsection*{Your turn}
\begin{itemize}
\raggedright
  \item \emph{Say it:} \emph{ārtī}
  \item \emph{Say it:} \emph{ārtī}, once more
  \item \emph{Say it:} \emph{sā̃je mandirmā̃ ārtī thaī.} — In the evening the aarti was held at the temple.
  \item \emph{Recall:} say the Gujarati for a holy day, then the Gujarati for worship, then say \emph{ārtī} again
\end{itemize}

\begin{tcolorbox}[breakable,colback=teal!4,colframe=teal!35!black,title={Before you move on}]
What does \gu{આરતી} mean? (\textquotedblleft{}An aarti (worship with lamps)\textquotedblright{}.) Say it once more.
\end{tcolorbox}
\section[\texorpdfstring{\gu{પ્રસાદ} (prasād)}{પ્રસાદ (prasād)}]{\gu{પ્રસાદ} (prasād) — a blessed offering (of food)}

\label{lesson:GU-C328-prasad}



\begin{quote}
Before the new one: say the Gujarati for worship, then the Gujarati for an aarti.
\end{quote}

\subsection*{You'll want to know: \gu{પ્રસાદ}}
\textbf{\gu{પ્રસાદ}} — \emph{prasād} — \textquotedblleft{}a blessed offering (of food)\textquotedblright{}.

\begin{quote}
\gu{આરતી} \gu{પછી} \gu{દાદીએ} \gu{બધાને} \gu{પ્રસાદ} \gu{આપ્યો}.
\end{quote}

\begin{quote}
\emph{ārtī pachī dādīe badhāne prasād āpyo.}
\end{quote}

\begin{quote}
After the aarti my grandmother gave everyone some of the blessed food.
\end{quote}

A masculine noun.

\begin{cousinweb}
\textbf{\gu{પ્રસાદ}} (\emph{prasād}) is a tatsama, Sanskrit \emph{prasāda}, \textquotedblleft{}clearness, kindness, favour\textquotedblright{} (Monier-Williams): the sweet or the fruit offered to a god and then shared out as the god's favour.
\end{cousinweb}

\begin{grammarlens}[title={What you've built}]
One more word for telling what happened.
\end{grammarlens}

\subsection*{Your turn}
\begin{itemize}
\raggedright
  \item \emph{Say it:} \emph{prasād}
  \item \emph{Say it:} \emph{prasād}, once more
  \item \emph{Say it:} \emph{ārtī pachī dādīe badhāne prasād āpyo.} — After the aarti my grandmother gave everyone some of the blessed food.
  \item \emph{Recall:} say the Gujarati for worship, then the Gujarati for an aarti, then say \emph{prasād} again
\end{itemize}

\begin{tcolorbox}[breakable,colback=teal!4,colframe=teal!35!black,title={Before you move on}]
What does \gu{પ્રસાદ} mean? (\textquotedblleft{}A blessed offering (of food)\textquotedblright{}.) Say it once more.
\end{tcolorbox}
\section[\texorpdfstring{\gu{પ્રગટાવવું} (pragṭāvvũ)}{પ્રગટાવવું (pragṭāvvũ)}]{\gu{પ્રગટાવવું} (pragṭāvvũ) — to light (a lamp)}

\label{lesson:GU-C328-pragtavvu}



\begin{quote}
Before the new one: say the Gujarati for an aarti, then the Gujarati for a blessed offering.
\end{quote}

\subsection*{You'll want to know: \gu{પ્રગટાવવું}}
\textbf{\gu{પ્રગટાવવું}} — \emph{pragṭāvvũ} — \textquotedblleft{}to light (a lamp)\textquotedblright{}.

\begin{quote}
\gu{મીરાએ} \gu{સાંજે} \gu{દીવો} \gu{પ્રગટાવ્યો}.
\end{quote}

\begin{quote}
\emph{Mīrāe sā̃je dīvo pragṭāvyo.}
\end{quote}

\begin{quote}
Mira lit the lamp in the evening.
\end{quote}

Past, with \gu{તેણે} (he or she …): \gu{તેણે} \gu{દીવો} \gu{પ્રગટાવ્યો}, \emph{teṇe dīvo pragṭāvyo}.

\begin{cousinweb}
\textbf{\gu{પ્રગટાવવું}} (\emph{pragṭāvvũ}) is built on Sanskrit \emph{prakaṭa}, \textquotedblleft{}manifest, visible\textquotedblright{} (CDIAL 8428): to make a flame appear. A lamp or a candle \textbf{\gu{પ્રગટાવે} \gu{છે}}; \textbf{\gu{સળગાવવું}} is for fires.
\end{cousinweb}

\begin{grammarlens}[title={What you've built}]
One more word for telling what happened.
\end{grammarlens}

\subsection*{Your turn}
\begin{itemize}
\raggedright
  \item \emph{Say it:} \emph{pragṭāvvũ}
  \item \emph{Say it:} \emph{pragṭāvvũ}, once more
  \item \emph{Say it:} \emph{Mīrāe sā̃je dīvo pragṭāvyo.} — Mira lit the lamp in the evening.
  \item \emph{Recall:} say the Gujarati for an aarti, then the Gujarati for a blessed offering, then say \emph{pragṭāvvũ} again
\end{itemize}

\begin{tcolorbox}[breakable,colback=teal!4,colframe=teal!35!black,title={Before you move on}]
What does \gu{પ્રગટાવવું} mean? (\textquotedblleft{}To light (a lamp)\textquotedblright{}.) Say it once more.
\end{tcolorbox}
\section[\texorpdfstring{\gu{ધૂપ} (dhūp)}{ધૂપ (dhūp)}]{\gu{ધૂપ} (dhūp) — incense}

\label{lesson:GU-C328-dhup}



\begin{quote}
Before the new one: say the Gujarati for a blessed offering, then the Gujarati for to light.
\end{quote}

\subsection*{You'll want to know: \gu{ધૂપ}}
\textbf{\gu{ધૂપ}} — \emph{dhūp} — \textquotedblleft{}incense\textquotedblright{}.

\begin{quote}
\gu{ઓરડામાં} \gu{ધૂપનો} \gu{ધુમાડો} \gu{હતો}.
\end{quote}

\begin{quote}
\emph{orḍāmā̃ dhūpno dhumāḍo hato.}
\end{quote}

\begin{quote}
There was incense smoke in the room.
\end{quote}

A masculine noun.

\begin{cousinweb}
\textbf{\gu{ધૂપ}} (\emph{dhūp}) is a tatsama, Sanskrit \emph{dhūpa}, \textquotedblleft{}incense, aromatic smoke\textquotedblright{} (Monier-Williams). It is masculine: \textbf{\gu{ધૂપનો} \gu{ધુમાડો}}, the smoke of incense.
\end{cousinweb}

\begin{grammarlens}[title={What you've built}]
One more word for telling what happened.
\end{grammarlens}

\subsection*{Your turn}
\begin{itemize}
\raggedright
  \item \emph{Say it:} \emph{dhūp}
  \item \emph{Say it:} \emph{dhūp}, once more
  \item \emph{Say it:} \emph{orḍāmā̃ dhūpno dhumāḍo hato.} — There was incense smoke in the room.
  \item \emph{Recall:} say the Gujarati for a blessed offering, then the Gujarati for to light, then say \emph{dhūp} again
\end{itemize}

\begin{tcolorbox}[breakable,colback=teal!4,colframe=teal!35!black,title={Before you move on}]
What does \gu{ધૂપ} mean? (\textquotedblleft{}Incense\textquotedblright{}.) Say it once more.
\end{tcolorbox}
